\chapter{La desnaturalización de la responsabilidad extracontractual del Estado en la Ley 2195 de 2022}
\section{Introducción: el flagelo de la corrupción}
\begin{itemize}
\item Uno de los principales males que acecha a Colombia.
\item El impacto de la corrupción en el patrimonio público.
\item Necesidad de instrumentos jurídicos para combatirla.
\item ¿Excesivo fetichismo legal?
\end{itemize}
\section{La reparación directa antes y después de la Ley 2195}
\subsection{Art. 140 de la Ley 1437 de 2011}
\textbf{Finalidad:} reparación del daño antijurídico.

{\small\begin{xltabular}{\linewidth}{|X|X|}\hline
\rowcolor{navy!12}\textbf{Demandante} & \textbf{Demandado} \\\hline\endhead
Persona interesada & Entidad pública causante del daño \\\hline
Entidad pública perjudicada & Particular / entidad causante del daño \\\hline
Persona interesada & Entidad pública por hechos de particular que haya obrado siguiendo una expresa instrucción de una entidad pública \\\hline
\end{xltabular}}
\subsection{Reformas del medio de control de reparación directa en la Ley 2195 (arts. 59 y 61)}
{\small\begin{xltabular}{\linewidth}{|X|X|}\hline
\rowcolor{navy!12}\textbf{Demandante} & \textbf{Demandado} \\\hline\endhead
Entidad pública lesionada / Agencia Nacional de Defensa Jurídica del Estado & Particular que ejerce función administrativa / servidor público causante del daño \\\hline
\end{xltabular}}
\begin{itemize}
\item \textbf{Finalidad:} reparación del daño al patrimonio público causado por un acto de corrupción (responsabilidad extracontractual derivada del delito).
\item Imposición de multas al responsable del acto de corrupción en la sentencia declaratoria de responsabilidad, por hasta 1.000 SMLMV, según la gravedad de la conducta.
\end{itemize}
\section{Responsabilidad extracontractual derivada del delito}
\begin{cita}«Artículo 1494. Las obligaciones nacen, ya del concurso real de las voluntades de dos o más personas, como en los contratos o convenciones; ya de un hecho voluntario de la persona que se obliga, como en la aceptación de una herencia o legado y en todos los cuasicontratos; ya a consecuencia de un hecho que ha inferido injuria o daño a otra persona, como en los delitos; ya por disposición de la ley, como entre los padres y los hijos de familia.» (Código Civil)\end{cita}
\section{Exposición de motivos: la supuesta insuficiencia de los mecanismos existentes}
\begin{cita}«Aunque el incidente de reparación integral y el procedimiento administrativo de responsabilidad fiscal tienen fines reparatorios, al igual que la acción de reparación directa que se propone, lo cierto es que no son suficientes para los propósitos del proyecto de ley por varias razones: i) existe una limitación en el reclamo de perjuicios, estos procesos se circunscriben exclusivamente a la solicitud de perjuicios materiales derivados del detrimento patrimonial, lo cual no permitiría reclamar otro tipo de perjuicios, por ejemplo, no se podría reclamar el daño al buen nombre "good will" del Estado; ii) no es claro que se puedan solicitar medidas de reparación no pecuniarias, como de satisfacción, por ejemplo unas excusas públicas o que se realicen medidas por parte del responsable para renovar la confianza entre la sociedad y la administración pública; iii) la responsabilidad fiscal es exclusiva de los servidores públicos y particulares que ejerzan función fiscal; iv) para que se dé el incidente de reparación en materia penal es necesario que se pruebe la responsabilidad penal del investigado; la acción de reparación directa propuesta solo exige que se pruebe el daño y que este haya ocurrido con ocasión de un acto de corrupción.»\end{cita}
\section{Reparación del daño e insuficiencia de los instrumentos existentes}
\begin{cita}«El daño al patrimonio público puede ser resarcido a través de medidas de reparación pecuniarias y no pecuniarias.»\end{cita}
\begin{itemize}
\item Art. 94 de la Ley 599 de 2000: «La conducta punible origina obligación de reparar los daños materiales y morales causados con ocasión de aquella». Sentencia C-344 de 2017: «en el entendido de que las categorías de perjuicios allí indicadas (…) no excluyen la reparación integral de todos los perjuicios tanto materiales, como inmateriales, que hayan sido causados a las víctimas como consecuencia del delito y resulten debidamente probados».
\item CSJ, Sala Penal, sentencia del 3 de mayo de 2017, Exp. 36784, M. P. Fernando Alberto Castro Caballero.
\item Medidas no pecuniarias: ¿era necesario crear una nueva acción para este tipo de medidas?
\end{itemize}
\section{Reparación de perjuicios}
\begin{cita}«El juez deberá tener en cuenta para la tasación de los perjuicios el impacto en la sociedad del acto de corrupción.»\end{cita}
\begin{itemize}
\item Función de la responsabilidad extracontractual: reparar el daño, todo el daño y nada más que el daño.
\item ¿Reducción de la indemnización por impacto a la sociedad?
\item ¿Incremento de la indemnización (daños punitivos) por impacto a la sociedad?
\end{itemize}
\section{¿Qué es un acto de corrupción y cómo se determina su ocurrencia?}
\begin{cita}«PARÁGRAFO 1o. Entiéndase por acto de corrupción las conductas penales enlistadas en los capítulos de delitos contra la administración pública, el medio ambiente, el orden económico y social, financiación del terrorismo y de grupos de delincuencia organizada, administración de recursos relacionados con actividades terroristas y de la delincuencia organizada, los consagrados en la Ley 1474 de 2011, los delitos electorales o cualquier conducta punible relacionada con el patrimonio público, que hubieren sido realizados.»\end{cita}
\begin{itemize}
\item ¿Debe hacer el juez de la reparación un juicio de adecuación típica? ¿Además debe estudiar la antijuridicidad y la culpabilidad de la conducta?
\item ¿Existe prejudicialidad respecto del proceso penal?
\end{itemize}
\section{\emph{Non bis in idem}}
\begin{cita}«PARÁGRAFO 2o. El pago que haya realizado el demandado en desarrollo de otro proceso judicial o fiscal de responsabilidad por los hechos de corrupción objeto del medio de control de reparación directa, se descontará del monto de la condena del proceso de reparación directa. De igual manera, en los otros procesos de responsabilidad en los cuales el demandado deba realizar un pago por el daño causado al patrimonio público, se descontará la suma reconocida y pagada en la sentencia de reparación directa.»\end{cita}
\begin{itemize}
\item ¿Realmente era necesario crear un nuevo instrumento procesal para ejercer pretensiones que se pueden perseguir a través de otros?
\end{itemize}
\section{Caducidad}
\begin{cita}«PARÁGRAFO 3o. El término para formular la pretensión de reparación directa derivada de un acto de corrupción se contará a partir del día siguiente de la fecha en que la entidad pública afectada tuvo o debió tener conocimiento de este, o en su defecto desde la ejecutoria del fallo definitivo adoptado en el proceso penal.»\end{cita}
\begin{itemize}
\item Dos extremos iniciales concurrentes para contar el término de caducidad (subjetivo / objetivo).
\item ¿Aplicación de la teoría del conocimiento del daño para personas jurídicas?
\end{itemize}
Puede representarse con dos líneas de tiempo:

{\small\begin{xltabular}{\linewidth}{|X|X|X|}\hline
\rowcolor{navy!12}\textbf{Línea de tiempo} & \textbf{Extremo inicial} & \textbf{Extremo final} \\\hline\endhead
Primera & Acción u omisión causante del daño & Conocimiento del daño \\\hline
Segunda & Conocimiento del daño & Ejecutoria del fallo penal \\\hline
\end{xltabular}}
\section{Imposición de multas (art. 61)}
\begin{cita}«PARÁGRAFO. Cuando la sentencia sea declaratoria de responsabilidad en los medios de control de reparación directa y controversias contractuales y el daño haya sido causado por un acto de corrupción, el juez deberá imponer, adicional al daño probado en el proceso, multa al responsable hasta de mil (1.000) salarios mínimos mensuales legales vigentes, la cual atenderá a la gravedad de la conducta, el grado de participación del demandado y su capacidad económica. El pago de la multa impuesta deberá dirigirse al Fondo de Reparación de las Víctimas de Actos de Corrupción. En la sentencia se deberán decretar las medidas cautelares que garanticen el pago de la sanción.»\end{cita}
\begin{itemize}
\item ¿Desnaturalización del medio de control de reparación directa?
\item ¿Vinculación del agente responsable y garantía del debido proceso?
\item ¿Otros medios de control distintos a la reparación directa y a las controversias contractuales?
\end{itemize}
\section{Problemas prácticos del \emph{non bis in idem} entre la reparación directa y el proceso de responsabilidad fiscal}
\begin{itemize}
\item ¿La sentencia absolutoria en el proceso de reparación directa hace tránsito a cosa juzgada en el proceso de responsabilidad fiscal? Depende del fundamento de la absolución: el demandado no lo cometió / inexistencia del hecho / inexistencia del acto de corrupción (problema de tipicidad).
\item ¿Puede existir diferencia en la tasación de perjuicios entre el proceso de reparación directa y el proceso de responsabilidad fiscal?
\end{itemize}
{\small\begin{xltabular}{\linewidth}{|X|X|}\hline
\rowcolor{navy!12}\textbf{Proceso} & \textbf{Objeto} \\\hline\endhead
Reparación directa & Reparación integral del daño al patrimonio público causado por acto de corrupción (art. 59 de la Ley 2195) \\\hline
Proceso de responsabilidad fiscal & Resarcimiento del daño ocasionado al patrimonio público como consecuencia de la conducta dolosa o culposa de quien realiza gestión fiscal, mediante el pago de una indemnización pecuniaria que compense el perjuicio sufrido por la entidad estatal (art. 4 de la Ley 610) \\\hline
\end{xltabular}}
\section{Función de la responsabilidad extracontractual del Estado}
{\small\begin{xltabular}{\linewidth}{|X|X|}\hline
\rowcolor{navy!12}\textbf{Fundamento} & \textbf{Función} \\\hline\endhead
Falla del servicio & Función demarcatoria \\\hline
Daño antijurídico & Función indemnizatoria \\\hline
Incremento de perjuicios por acto de corrupción & Función sancionatoria \\\hline
\end{xltabular}}
\section{Conclusiones}
\begin{itemize}
\item El populismo punitivo en la lucha contra la corrupción.
\item Suficiencia de los mecanismos existentes.
\item Inconsistencia de las modificaciones al medio de control de reparación directa creadas en la Ley 2195.
\item Desnaturalización del juicio de responsabilidad extracontractual / confusión con el juicio de responsabilidad penal.
\end{itemize}
